% ---- pg-523 （书页 511） ----

\noindent\hfill{\zihao{6}续表}\par\vspace{2pt}

\begin{mvtable}
\begin{longtable}{cccccc}\toprule
行 & $X_1=\dfrac{\mathrm{CF}}{\mathrm{TD}}$ & $X_2=\dfrac{\mathrm{NI}}{\mathrm{TA}}$ & $X_3=\dfrac{\mathrm{CA}}{\mathrm{CL}}$ & $X_4=\dfrac{\mathrm{CA}}{\mathrm{NS}}$ & 总体 $\pi_i,i=1,2$\\
\midrule
\endfirsthead
\multicolumn{6}{c}{\zihao{6}续表}\\
\toprule
行 & $X_1=\dfrac{\mathrm{CF}}{\mathrm{TD}}$ & $X_2=\dfrac{\mathrm{NI}}{\mathrm{TA}}$ & $X_3=\dfrac{\mathrm{CA}}{\mathrm{CL}}$ & $X_4=\dfrac{\mathrm{CA}}{\mathrm{NS}}$ & 总体 $\pi_i,i=1,2$\\
\midrule
\endhead
\bottomrule
\endlastfoot

19 & 0.01 & $-0.00$ & 1.26 & 0.60 & 0\\
20 & 0.12 & 0.11 & 1.14 & 0.17 & 0\\
21 & $-0.28$ & $-0.27$ & 1.27 & 0.51 & 0\\
1 & 0.51 & 0.10 & 2.49 & 0.54 & 1\\
2 & 0.08 & 0.02 & 2.01 & 0.53 & 1\\
3 & 0.38 & 0.11 & 3.27 & 0.35 & 1\\
4 & 0.19 & 0.05 & 2.25 & 0.33 & 1\\
5 & 0.32 & 0.07 & 4.24 & 0.63 & 1\\
6 & 0.31 & 0.05 & 4.45 & 0.69 & 1\\
7 & 0.12 & 0.05 & 2.52 & 0.69 & 1\\
8 & $-0.02$ & 0.02 & 2.05 & 0.35 & 1\\
9 & 0.22 & 0.08 & 2.35 & 0.40 & 1\\
10 & 0.17 & 0.07 & 1.80 & 0.52 & 1\\
11 & 0.15 & 0.05 & 2.17 & 0.55 & 1\\
12 & $-0.10$ & $-0.01$ & 2.50 & 0.58 & 1\\
13 & 0.14 & $-0.03$ & 0.46 & 0.26 & 1\\
14 & 0.14 & 0.07 & 2.61 & 0.52 & 1\\
15 & 0.15 & 0.06 & 2.23 & 0.56 & 1\\
16 & 0.16 & 0.05 & 2.31 & 0.20 & 1\\
17 & 0.29 & 0.06 & 1.84 & 0.38 & 1\\
18 & 0.54 & 0.11 & 2.33 & 0.48 & 1\\
19 & $-0.33$ & $-0.09$ & 3.01 & 0.47 & 1\\
20 & 0.48 & 0.09 & 1.24 & 0.18 & 1\\
21 & 0.56 & 0.11 & 4.29 & 0.45 & 1\\
22 & 0.20 & 0.08 & 1.99 & 0.30 & 1\\
23 & 0.47 & 0.14 & 2.92 & 0.45 & 1\\
24 & 0.17 & 0.04 & 2.45 & 0.14 & 1\\
25 & 0.58 & 0.04 & 5.06 & 0.13 & 1\\
\end{longtable}
\end{mvtable}

\noindent{\zihao{6}注: $\pi_1=0$: 破产企业; $\pi_2=1$: 非破产企业.}

\noindent{\zihao{6}资料来源: 1968,1969,1970,1971,1972, 穆迪行业手册(Moody's Industrial Manuals).}

\noindent (a)\quad 对每一组数据用不同的符号, 对以下各对变量的观察值作平面散布图: $(x_1,x_2)$, $(x_1,x_3)$ 及 $(x_1,x_4)$. 是否任何一对变量的数据散布图都显示二元正态性?

\noindent (b)\quad 用破产企业的 $n_1=21$ 对 $(x_1,x_2)$ 观测值及非破产企业的 $n_2=25$ 对 $(x_1,x_2)$ 观测值分别计算样本均值向量 $\bar{x}_1$ 和 $\bar{x}_2$, 以及样本协方差矩阵 $\mathbf{S}_1$ 和 $\mathbf{S}_2$.

\noindent (c)\quad 利用(b)中的结果并假定两组随机样本都服从二元正态分布, 构造在 $p_1=p_2$ 及 $c(1|2)=c(2|1)$ 时的分类法则(11-29).

\noindent (d)\quad 对(c)中的分类法则的表现进行评估, 评估时利用两种失误率: \hyperref[eq:11-34]{式(11-34)}中的表观失误率(APER)及\hyperref[eq:11-36]{式(11-36)}中估计的期望真实失误率 $\hat{E}(\mathrm{AER})$.
